% Generated draft: requires mathematical and source review.
\documentclass[11pt]{article}
\usepackage[margin=25mm]{geometry}
\usepackage{fontspec}
\setmainfont{Latin Modern Roman}
\usepackage{amsmath,amssymb,mathtools,unicode-math}
\setmathfont{Latin Modern Math}
\usepackage{enumitem,needspace,xurl,hyperref}
\setlength{\parindent}{0pt}
\setlength{\parskip}{5pt}
\setlength{\emergencystretch}{3em}
\newcommand{\N}{\mathbb{N}}
\newcommand{\Z}{\mathbb{Z}}
\newcommand{\Q}{\mathbb{Q}}
\newcommand{\R}{\mathbb{R}}
\newcommand{\C}{\mathbb{C}}
\newcommand{\F}{\mathbb{F}}
\newcommand{\A}{\mathbb{A}}
\newcommand{\PP}{\mathbb{P}}
\newcommand{\E}{\mathbb{E}}
\newcommand{\eps}{\varepsilon}
\newcommand{\dd}{\,\mathrm{d}}
\begin{document}
\section*{UnsolvedMath ID: 30004737}
\textbf{Source problem code:} catalogue entry
\subsection{Definitions and mathematical statement}
The intended question concerns an affine map \(a\) described in the supplied record as Diophantine and step-two. The record does not specify the underlying space, the definition of step-two, the Diophantine condition, or the meaning and required smallness of a perturbation. These qualifications remain unresolved; the following expresses the conjugacy question conditionally on their specification.

Let \(M\) denote the underlying smooth manifold, which must be specified, and write \(\operatorname{Diff}^{r}(M)\) for the diffeomorphisms of \(M\) whose maps and inverses are of class \(C^{r}\). Interpret smooth as \(C^{\infty}\). For every admissible Diophantine step-two affine diffeomorphism \(a\) of \(M\), and every \(f\in\operatorname{Diff}^{\infty}(M)\) satisfying the intended perturbation hypothesis relative to \(a\), is the following implication true?
\[
\left(\exists h\in\operatorname{Diff}^{1}(M)\;:\;h\circ f=a\circ h\right)
\quad\Longrightarrow\quad
\left(\exists H\in\operatorname{Diff}^{\infty}(M)\;:\;H\circ f=a\circ H\right).
\]
The conclusion asks for the existence of a smooth conjugacy; it does not require the particular given conjugacy \(h\) to be smooth. This is a conditional formulation of the supplied question, not a fully specified assertion until the missing geometric, arithmetic, and perturbation hypotheses are supplied.
\subsection{Short English statement}
Must every smooth perturbation of a Diophantine step-two affine map that admits a \(C^{1}\) conjugacy also admit a smooth conjugacy? The underlying space, arithmetic condition, step-two condition, and perturbation hypothesis are unspecified, so the formulation requires human mathematical review.
\subsection{Sources}
\begin{itemize}
\item Supplied local UnsolvedMath registry record, permanent ID \(30004737\), exact source problem code \texttt{catalogue entry}, title ``Smooth Rigidity of Diophantine Step-Two Affine Maps,'' in the Oberwolfach Reports Open Problems collection. The record gives proposed year \(2021\) and names no proposer.
\item Supplied report citation: \emph{Dynamische Systeme} (2022), Oberwolfach Reports, DOI 10.4171/owr/2021/33, pp.\ 1735--1803. \url{https://doi.org/10.4171/owr/2021/33}. The supplied background describes these reports as documenting mathematical workshops and their open-problem discussions. The report contents have not been independently consulted here.
\item Supplied literature reference, labeled ``parabolic KAM rigidity'': \url{https://doi.org/10.1007/s00222-026-01415-7}.
\item Supplied literature reference, labeled ``higher-rank local rigidity'': \url{https://arxiv.org/abs/2106.01585}.
\item Literature assessment dated exactly \texttt{2026-08-22}: the supplied curation labels the problem partially solved and states that the specific one-map \(C^{1}\)-to-smooth conjugacy question is not settled in full generality. It further attributes to 2026 work a locked--unlocked dichotomy and tame KAM rigidity for almost every translation in parabolic \(\mathbb{Z}^{2}\) actions with a step-two generator, and states that higher-step cases remain open. This is supplied dated curation derived from the OWR collection, not an independently verified literature assessment. Its time-sensitive status and cited claims require independent verification; the higher-rank results are not asserted here to settle the one-map question.
\end{itemize}
\end{document}
